\section{Results}
\label{sec:results}

\subsection{C1: Behavioral Signals Shift Under Poisoning}
\label{sec:results:c1}

\begin{figure}[tbp]
  \centering
  \figureiffound{../results/signals_minja.png}%
    {Phase 1 \texttt{signals\_minja} panel; produced by the signal
    validation run}
  \caption{Phase 1 signal validation: the six behavioral signals per
  PR-handling step, 25 clean steps then 25 poisoned steps, under
  MINJA.}
  \label{fig:signals-minja}
\end{figure}

\Cref{fig:signals-minja} reports the per-step values of the six
behavioral signals. KL divergence shifts at the injection boundary,
with Cohen's $d = 1.343$ under MINJA and $d = 1.402$ under AgentPoison
(\texttt{results/decision\_matrix.csv}).

\subsection{C2: Advance Warning Time}
\label{sec:results:c2}

\begin{figure}[tbp]
  \centering
  \figureiffound{../results/main_experiment/figures/fig1_signal_drift.png}%
    {Per-experiment signal-drift panel for \texttt{main\_experiment};
    produced by \texttt{plot\_signal\_drift}.}
  \caption{Signal-drift trajectories for the \texttt{main\_experiment}
  run; shaded band marks the post-injection window.}
  \label{fig:signal-drift}
\end{figure}

\Cref{fig:signal-drift} plots every behavioral signal across the full
PR stream for a single \texttt{main\_experiment} run, and
\cref{tab:main-results} gives the per-experiment \awt{}, AUROC, F1,
and allocation efficiency.

\begin{table}[tbp]
  \centering
  \caption{Headline results: per-experiment \awt{}, AUROC, F1, and
  allocation efficiency over the Phase 10 run set. Generated by
  \texttt{make\_main\_results\_table}.}
  \label{tab:main-results}
  \inputiffound{../results/tables/table1_main_results.tex}%
    {headline AWT / AUROC / efficiency table, generated by
    \texttt{make\_main\_results\_table}.}
\end{table}

\subsection{C3: Allocation Efficiency Under Attack}
\label{sec:results:c3}

\begin{figure}[tbp]
  \centering
  \figureiffound{../results/figures/fig4_allocation_efficiency_over_time.png}%
    {Multi-experiment allocation efficiency comparison; produced by
    \texttt{plot\_allocation\_efficiency\_over\_time}.}
  \caption{Throughput per step, \sys{} vs.\ no-monitoring and vs.\
  \texttt{ablation\_no\_health\_scores}.}
  \label{fig:alloc-eff}
\end{figure}

\Cref{fig:alloc-eff} shows the cumulative successful-task rate per
step for three configurations, and \cref{tab:ablation}, generated by
\texttt{make\_ablation\_table}, gives the numbers behind the
comparison.

\subsection{C4: AUROC vs.\ Sentinel}
\label{sec:results:c4}

\begin{figure}[tbp]
  \centering
  \figureiffound{../results/figures/fig5_roc_curve.png}%
    {Multi-experiment ROC comparison; produced by
    \texttt{plot\_roc\_curve}.}
  \caption{ROC curves for \sys{} vs.\ \texttt{baseline\_sentinel} on
  both attacks.}
  \label{fig:roc}
\end{figure}

\Cref{fig:roc} overlays the ROC curves for \sys{} and
\texttt{baseline\_sentinel} across both attack types;
\cref{tab:main-results} lists the AUROC of each.

\subsection{C5: Forecaster Overhead}
\label{sec:results:c5}

\Cref{tab:latency} compares the mean per-step wall-clock latency
across detector configurations, BOCPD alone against the full
BOCPD + Chronos T5-small ensemble.

\begin{table}[tbp]
  \centering
  \caption{Per-step detector latency (ms). Lower is better. Generated
  by \texttt{make\_latency\_table}.}
  \label{tab:latency}
  \inputiffound{../results/tables/table4_latency.tex}%
    {latency comparison table, generated by
    \texttt{make\_latency\_table}.}
\end{table}

\subsection{Ablations}
\label{sec:results:ablations}

\begin{table}[tbp]
  \centering
  \caption{Per-knob ablation deltas relative to the full system.
  Generated by \texttt{make\_ablation\_table}.}
  \label{tab:ablation}
  \inputiffound{../results/tables/table2_ablation.tex}%
    {ablation deltas table, generated by \texttt{make\_ablation\_table}.}
\end{table}

\begin{figure}[tbp]
  \centering
  \figureiffound{../results/figures/fig6_ablation_bar_chart.png}%
    {Multi-experiment ablation bar chart; produced by
    \texttt{plot\_ablation\_bar\_chart}.}
  \caption{Per-ablation Advance Warning Time relative to the full
  system.}
  \label{fig:ablation-bar}
\end{figure}

\Cref{fig:ablation-bar} renders the per-knob \awt{} delta relative to
the full system, complementing the per-metric numbers in
\cref{tab:ablation}. The forecaster channel exercised by the experiment
runner is an exponential-moving-average residual stand-in for the
Chronos model (\texttt{full\_system\_detector.py},
\texttt{FORECASTER\_IS\_PLACEHOLDER}), so the forecaster ablation
isolates the contribution of that stand-in; re-running it through
\texttt{ChronosForecaster} is the next ablation run.
